\documentclass[tikz,border=8pt]{standalone}
\usepackage{tikz}
\usepackage{amsmath}
\usepackage{pgfplots}
\pgfplotsset{compat=1.18}
\usepackage{ctex}
\usetikzlibrary{calc, positioning}

\begin{document}
\begin{tikzpicture}

% 左子图：Kruskal-Wallis 3组箱线图
\begin{axis}[
  name=ax1,
  width=7cm, height=6.5cm,
  xmin=0.3, xmax=3.7,
  ymin=0, ymax=13.5,
  xlabel={组别},
  ylabel={观测值},
  axis lines=left,
  xtick={1,2,3},
  xticklabels={组\,1,组\,2,组\,3},
  xticklabel style={font=\small},
  ytick={0,2,4,6,8,10,12},
  yticklabel style={font=\scriptsize},
  grid=major, grid style={gray!30},
  clip=false
]
  % 组1箱线图：IQR [2,5], 中位线 3.5, 须 [1,7]
  \addplot[fill=blue!20, draw=blue!70!black, thick]
    coordinates {(0.7,2)(1.3,2)(1.3,5)(0.7,5)(0.7,2)};
  \addplot[blue!70!black, ultra thick]
    coordinates {(0.7,3.5)(1.3,3.5)};
  \addplot[blue!70!black, thick] coordinates {(1,5)(1,7)};
  \addplot[blue!70!black, thick] coordinates {(1,2)(1,1)};
  \addplot[blue!70!black, thick] coordinates {(0.8,7)(1.2,7)};
  \addplot[blue!70!black, thick] coordinates {(0.8,1)(1.2,1)};

  % 组2箱线图：IQR [4,7], 中位线 5.5, 须 [3,9]
  \addplot[fill=red!20, draw=red!70!black, thick]
    coordinates {(1.7,4)(2.3,4)(2.3,7)(1.7,7)(1.7,4)};
  \addplot[red!70!black, ultra thick]
    coordinates {(1.7,5.5)(2.3,5.5)};
  \addplot[red!70!black, thick] coordinates {(2,7)(2,9)};
  \addplot[red!70!black, thick] coordinates {(2,4)(2,3)};
  \addplot[red!70!black, thick] coordinates {(1.8,9)(2.2,9)};
  \addplot[red!70!black, thick] coordinates {(1.8,3)(2.2,3)};

  % 组3箱线图：IQR [6,9], 中位线 7.5, 须 [5,11]
  \addplot[fill=green!20, draw=green!60!black, thick]
    coordinates {(2.7,6)(3.3,6)(3.3,9)(2.7,9)(2.7,6)};
  \addplot[green!60!black, ultra thick]
    coordinates {(2.7,7.5)(3.3,7.5)};
  \addplot[green!60!black, thick] coordinates {(3,9)(3,11)};
  \addplot[green!60!black, thick] coordinates {(3,6)(3,5)};
  \addplot[green!60!black, thick] coordinates {(2.8,11)(3.2,11)};
  \addplot[green!60!black, thick] coordinates {(2.8,5)(3.2,5)};

  % 合并秩说明
  \node[font=\scriptsize, align=left, fill=white, draw=gray!50, thin,
        rounded corners=2pt, inner sep=3pt]
    at (axis cs:2.0, 1.8) {合并排秩\\计算 $H$ 统计量};

\end{axis}

% 右子图：卡方拟合优度（观察 vs 期望柱状图）
\begin{axis}[
  name=ax2,
  at={(ax1.east)}, anchor=west, xshift=0.8cm,
  width=7cm, height=6.5cm,
  xmin=0.3, xmax=5.7,
  ymin=0, ymax=42,
  xlabel={类别},
  ylabel={频数},
  axis lines=left,
  xtick={1,2,3,4,5},
  xticklabels={A,B,C,D,E},
  xticklabel style={font=\small},
  ytick={0,5,10,15,20,25,30,35,40},
  yticklabel style={font=\scriptsize},
  grid=major, grid style={gray!30},
  clip=false,
  legend style={at={(1.02,0.5)}, anchor=west, font=\scriptsize,
                fill=white, draw=gray!50}
]
  % 观察频数 O（蓝色，左偏）
  \addplot+[ybar, bar width=0.25, xshift=-7pt,
            fill=blue!30, draw=blue!70!black]
    coordinates {(1,28)(2,22)(3,15)(4,18)(5,32)};
  \addlegendentry{观察频数 $O$}

  % 期望频数 E（橙色，右偏）
  \addplot+[ybar, bar width=0.25, xshift=7pt,
            fill=orange!30, draw=orange!80!black]
    coordinates {(1,23)(2,23)(3,23)(4,23)(5,23)};
  \addlegendentry{期望频数 $E$}

  % 卡方公式（左上角，不与图例重叠）
  \node[font=\small, align=center, fill=white, draw=gray!50, thin,
        rounded corners=2pt, inner sep=4pt]
    at (axis cs:1.8, 38) {$\chi^2=\displaystyle\sum\frac{(O-E)^2}{E}$};

\end{axis}

% 左子图 H 统计量公式（ax1 正下方，分两行显示）
\node[font=\small, align=center, fill=white, draw=gray!50, thin,
      rounded corners=2pt, inner sep=5pt]
  at ($(ax1.south)+(0,-1.4cm)$)
  {$H=\dfrac{12}{N(N{+}1)}\displaystyle\sum_{i=1}^{k}\dfrac{R_i^2}{n_i}-3(N{+}1)$\\[4pt]
   $H\sim\chi^2(k{-}1)$（$H_0$ 下）};

% 右子图卡方检验说明（ax2 正下方）
\node[font=\small, align=center, fill=white, draw=gray!50, thin,
      rounded corners=2pt, inner sep=5pt]
  at ($(ax2.south)+(0,-1.4cm)$)
  {$H_0$：数据服从指定分布\\[4pt]
   拒绝域：$\chi^2>\chi^2_\alpha(k{-}1{-}r)$};

% 子图标题（外置 node）
\node[font=\bfseries\small] at ($(ax1.north)+(0,0.4cm)$) {Kruskal-Wallis 检验（3 组）};
\node[font=\bfseries\small] at ($(ax2.north)+(0,0.4cm)$) {卡方拟合优度检验};

% 整体标题
\node[font=\large\bfseries] at ($(ax1.north)!0.5!(ax2.north)+(0,1.3cm)$)
  {Kruskal-Wallis 检验 \& 卡方拟合优度检验};

\end{tikzpicture}
\end{document}
